We compare HStore, a persistent hypergraph storage engine built on copy-on-write counted B+trees, with HyperGraphDB~1.4 on Berkeley DB Java Edition. Both systems run 26 storage-level workloads on the same skewed synthetic hypergraph, with up to 1.45 million incidences and equal 64\,MiB caches, in three configurations: two data sizes and two commit modes. Every read result is cross-checked between the systems. At scale 1 with asynchronous commits, HStore enumerates incidence sets \M{ratioAreadincidence} faster, counts shared hyperedges \M{ratioAreadcomembership} faster, applies batched updates \M{ratioAwriteupdate} faster, and needs \M{ratioAdisk} less disk. These leads come from storing incidence sets as small sorted trees, and from commits that log only the keys they changed while a background checkpoint writes pages. When the data is four times larger than the cache, HStore keeps most of these leads but falls to \M{sumCreadmembers} of HyperGraphDB's speed on member enumeration, because its member lookups miss the node cache. With synchronous commits, both systems are bound by the cost of one forced write per commit. We also report latency percentiles, bytes written, heap use, space reclamation and crash recovery, the variation over runs, and a durability setting in HyperGraphDB that, left unchecked, makes its synchronous commits unsynced.
